\section*{Resumo}
\addcontentsline{toc}{section}{Resumo}

O MQTT é um dos protocolos de mensagens mais usados na Internet das Coisas, mas deixa a
segurança à implementação. Como todas as mensagens atravessam um \emph{broker}, o
TLS não basta: termina no próprio \emph{broker}, que vê todo o conteúdo, e não
impede a repetição de mensagens nem o abuso de uma chave comprometida. Este trabalho
apresenta uma revisão sistemática da literatura sobre mecanismos de segurança para
o MQTT, conduzida segundo um protocolo explícito (perguntas de investigação, PICOC,
\emph{strings} de pesquisa, critérios de inclusão e exclusão, e extração de
dados). As pesquisas no IEEE Xplore, no Scopus e no ScienceDirect devolveram
\nTriados{} registos únicos, triados por título e resumo; \nIncluidos{} estudos
primários (2015--2026) e cinco revisões relacionadas foram analisados em texto
integral, classificados numa taxonomia de oito categorias e comparados pelas
propriedades de segurança que garantem. A literatura trata sobretudo
confidencialidade, autenticação e integridade. A validade temporal, a deteção de
duplicados, a ligação ao tópico e o não repúdio existem, mas isolados em poucos
estudos, e a maioria das soluções confia no \emph{broker}. Daqui resulta a
proposta: propõe-se uma arquitetura em
camadas comutáveis, com defesa no transporte, um envelope fim-a-fim com AES-256-GCM
(ou ChaCha20-Poly1305) ou Ed25519 e limitação de tráfego antes de qualquer operação criptográfica, e
planeia-se a segunda parte do projeto: medir, num Raspberry Pi, o custo e os
ataques impedidos por cada camada, para escolher a configuração com a maior
proteção e o menor impacto no desempenho.

\vspace{0.6cm}
\noindent\textbf{Palavras-chave:} MQTT; Internet das Coisas; segurança fim-a-fim;
cifra autenticada; revisão sistemática da literatura; defesa em profundidade.
